\section{The Dust on the Forehead: Mortality as the Wages of Sin}

Every Lent begins with a gesture that leaves nothing to evasion. Before the altar, the priest dips his thumb into ashes burned from last year's palms, traces the sign of the Cross upon each bowed forehead, and speaks the sentence pronounced over the first man outside Eden: \emph{Memento, homo, quia pulvis es, et in pulverem reverteris}---``Remember, man, that thou art dust, and into dust thou shalt return.'' Yet the liturgy never speaks of dust without first speaking of mercy. Immediately after the ashes are imposed, the Introit of the Ash Wednesday Mass rises in the words of Wisdom: \emph{Misereris omnium, Domine, et nihil odisti eorum quae fecisti}---``Thou hast mercy upon all, O Lord, and hatest none of the things which thou hast made,'' overlooking the sins of men for the sake of repentance.\endnote{1962 \emph{Missale Romanum}, Feria IV Cinerum, p.~57 (imposition of ashes adapting Gen.\ 3:19; Introit from Wis.\ 11:24--25, 27 with Ps.\ 56:2). Throughout, English Scripture follows the Challoner Douay-Rheims Bible, paired with the Clementine Vulgate.} Between that dust and that mercy lies the whole Catholic doctrine of death.

The Church begins her reflection on mortality neither with Stoic resignation to nature's neutral cycle nor with dualist contempt for the flesh, but where revelation begins: with a creature made in the image of God for unbroken life, and with an intrusion that ought never to have entered the garden.

\subsection{God Made Not Death}

Sacred Scripture is unambiguous on the origin of mortality. On the first pages of Genesis, the Lord God forms man from the slime of the earth, breathes into his face the breath of life (Gen.\ 2:7), and guards His single precept in paradise with a solemn warning: ``But of the tree of knowledge of good and evil, thou shalt not eat. For in what day soever thou shalt eat of it, thou shalt die the death'' (Gen.\ 2:17; \emph{morte morieris}). Death enters the biblical narrative not as an ornament of creation, but as the penalty of disobedience. When the precept is broken, the sentence falls: ``In the sweat of thy face shalt thou eat bread till thou return to the earth out of which thou wast taken: for dust thou art, and into dust thou shalt return'' (Gen.\ 3:19).

Centuries later, the Book of Wisdom gathered the mind of Israel into a classic pillar of Christian anthropology:

\begin{quote}
For God made not death, neither hath he pleasure in the destruction of the living. For he created all things that they might be: and he made the nations of the earth for health: and there is no poison of destruction in them, nor kingdom of hell upon the earth. [\ldots] For God created man incorruptible, and to the image of his own likeness he made him. But by the envy of the devil, death came into the world. (Wis.\ 1:13--14; 2:23--24)
\end{quote}

St.~Paul draws the same line from Eden to Calvary in his Epistle to the Romans: ``Wherefore as by one man sin entered into this world and by sin death: and so death passed upon all men, in whom all have sinned'' (Rom.\ 5:12). A chapter later, he condenses the Fall and the Redemption into one stark antithesis: ``For the wages of sin is death. But the grace of God, life everlasting in Christ Jesus our Lord'' (Rom.\ 6:23; \emph{Stipendia enim peccati, mors}). The Latin \emph{stipendia} denotes a soldier's pay---what is earned under a master's service. Sin pays its servants in dissolution; eternal life is never a wage earned by fallen clay, but the unmerited gift of God in Christ.

In its Fifth Session (17 June 1546), the Council of Trent solemnly defined this teaching in its Decree on Original Sin. Canon 1 anathematizes anyone who denies that Adam, upon transgressing God's commandment in paradise, immediately lost the holiness and justice in which he had been constituted, incurred divine wrath and indignation, and consequently fell under death and captivity to the devil, who held the empire of death. Canon 2 anathematizes anyone who claims that Adam's transgression harmed himself alone, or that he passed on to the whole human race merely bodily death and bodily penalties rather than sin itself, which is the death of the soul.\endnote{Council of Trent, Session V (17 June 1546), Decree on Original Sin, canons 1--2 (DS 1511--1512), appealing to Rom.\ 5:12.} Bodily death is genuinely the penalty of Adam's sin passed to all his stock; yet behind the death of the body stands the deeper death of the soul severed from God.

\subsection{Original Justice and the Gift of Immortality}

Yet a classic difficulty arises immediately. Man's body is an animal frame composed of contrary elements, sustained by food and naturally liable to wear and dissolution. How, then, would Adam not have died had he kept the commandment?

Both St.~Augustine and St.~Thomas Aquinas answer with a luminous distinction. In the opening chapter of Book XIII of \emph{The City of God}, Augustine explains that God did not make Adam like the angels, who cannot die a bodily death even after sinning:

\begin{quote}
He had so made them, that if they discharged the obligations of obedience, an angelic immortality and a blessed eternity might ensue, without the intervention of death; but if they disobeyed, death should be visited on them with just sentence.\endnote{St.~Augustine, \emph{The City of God}, XIII.1 (Marcus Dods trans., vol.~1, 1871).}
\end{quote}

In original innocence, Adam did not yet possess \emph{non posse mori} (the impossibility of dying, proper to glorified life in heaven), but he truly possessed \emph{posse non mori}---the real power not to die while obedient to God. That power was a preternatural gift rooted in original justice, the harmonious grace holding reason subject to God, lower appetites subject to reason, and the mortal body subject to the soul.

St.~Thomas Aquinas unfolds the metaphysics of this gift in the \emph{Summa Theologiae}. A thing may be incorruptible by nature (like an angel), by the overflow of glory in the risen body, or ``on the part of its efficient cause''---and in this third sense Adam's body ``was indissoluble not by reason of any intrinsic vigor of immortality, but by reason of a supernatural force given by God to the soul, whereby it was enabled to preserve the body from all corruption so long as it remained itself subject to God.''\endnote{St.~Thomas Aquinas, \emph{Summa Theologiae}, I, q.~97, a.~1 (English Dominican trans.); on ordinary food and the Tree of Life sustaining that life until translation to glory, see I, q.~97, aa.~3--4.}

Even clearer is Aquinas's reply in the \emph{Prima Secundae} to the objection that death is simply natural to flesh and blood. A smith forging a knife chooses iron because it is hard and takes an edge; that iron also rusts is an unchosen defect of the material, which the craftsman would banish if he could. So too human nature requires a composite, sensitive body to serve a rational soul, while corruptibility clings to matter as a defect repugnant to the immortal form animating it. God, lord of every nature, ``in forming man supplied the defect of nature, and by the gift of original justice, gave the body a certain incorruptibility.''\endnote{Aquinas, \emph{Summa Theologiae}, I-II, q.~85, a.~6.} Adam's sin caused death by pulling away that sustaining gift, ``as by displacing a pillar a man moves accidentally the stone resting thereon.'' Once the pillar of original justice fell, the soul lost its supernatural hold over the body and the flesh fell under the gravity of its corruptible elements---a withdrawal of favor that bears the true character of punishment ordered by divine justice.\endnote{Aquinas, \emph{Summa Theologiae}, I-II, q.~85, a.~5 and ad 1.}

\subsection{The First Death and the Second}

Because sin broke both man's union with God and his harmony of spirit and clay, Augustine warns that Christian faith never speaks of death in one sense alone. In Book XIII, chapter 2 of \emph{The City of God}, he maps the distinction governing all the last things:

\begin{quote}
The death, then, of the soul takes place when God forsakes it, as the death of the body when the soul forsakes it. Therefore the death of both---that is, of the whole man---occurs when the soul, forsaken by God, forsakes the body. For, in this case, neither is God the life of the soul, nor the soul the life of the body. And this death of the whole man is followed by that which, on the authority of the divine oracles, we call the second death.\endnote{St.~Augustine, \emph{The City of God}, XIII.2 (Dods, vol.~1); cf.\ Apoc.\ 20:14; Matt.\ 10:28.}
\end{quote}

Bodily death separates soul from body; spiritual death separates the soul from God through mortal sin; and dying in that alienation opens into the second death, perpetual exile from God. Hence Augustine's sharp paradox: ``Of the first and bodily death, then, we may say that to the good it is good, and evil to the evil. But, doubtless, the second, as it happens to none of the good, so it can be good for none.''\endnote{St.~Augustine, \emph{The City of God}, XIII.2.}

Already the Old Testament glimpsed that dying in God's friendship is peace: ``The souls of the just are in the hand of God, and the torment of death shall not touch them\ldots{} they are in peace'' (Wis.\ 3:1--3). Yet until the Redeemer came, Adam's shadow lay unbroken over every cradle and every tomb: ``Death reigned from Adam unto Moses, even over them also who have not sinned, after the similitude of the transgression of Adam, who is a figure of him who was to come'' (Rom.\ 5:14). The ashes of Ash Wednesday confess that man was not made to rot, and that mortality bears the scar of an ancient rebellion---yet the priest traces those very ashes in the sign of the Cross.
